\subsection{Computing \texorpdfstring{$A(p)$}{A(p)} in practice}
\label{a:sec:practical}

In the absence of an elementary formula, evaluation must combine exact numerical
iteration with asymptotic approximation.  The two regimes identified in
\cref{a:sec:Abounds} indicate where to divide the calculation.

Across most of the interval, direct iteration is satisfactory.  The infinite
product \cref{a:eq:prodFormula} converges quickly, each partial product bounds
$A(p)$ from above, and a modest number of iterations gives many digits.  The
difficulty is concentrated near $p=\tfrac12$.  As $A(p)\to0$, the quantity of
interest $1/A$ diverges, so a fixed absolute error in $A$ becomes unbounded in the
quasi-stationary mean.  More importantly, the process lifetime grows sharply near
criticality and $S_n$ recedes ever more slowly.  Iteration must continue until the
ratio has effectively stabilised, at a cost illustrated by \cref{a:tab:Avals}:
about $10^2$ iterations at $p=0.2$, $10^3$ at $p=0.45$, and nearly
$5\times10^5$ at $p=0.4999$.

The near-critical region can instead be handled by the asymptotic
\cref{a:eq:A-near-critical}, which is exact in the limit and inexpensive to
evaluate:
\begin{equation}
\label{a:eq:Apractical}
A(p) \approx
\begin{cases}
\displaystyle \frac12\prod_{k=1}^{N}\lrb{1-\frac{S_k}{2}}, & p<p_\ast,\quad N\gg1,\\[14pt]
2-4p, & p\ge p_\ast,
\end{cases}
\end{equation}
with the crossover $p_\ast$ chosen just below $\tfrac12$.
\Cref{a:thm:A-near-critical} indicates where that crossover should lie.  The
relative error of the asymptotic is about
$\varepsilon\ln(1/\varepsilon)$, so at $p_\ast = 0.4999$
($\varepsilon = 2\times10^{-4}$) the asymptotic is already accurate to about two
parts in $10^3$, while the product still converges in a few hundred thousand
iterations.  Any $p_\ast$ between $0.499$ and $0.4999$ gives better than one per
cent accuracy throughout.  If continuity of the derivative is required, the two
branches may be blended across a narrow window.

This hybrid calculation produces the numerical values quoted in this chapter and
in \cref{a:tab:Avals}; the corresponding conditional-mean plot appears in \ChM.

\FloatBarrier
